\chapter{Conclusiones y trabajos futuros}

\section{Conclusiones}

El trabajo ha permitido construir un flujo reproducible para analizar biomarcadores transcriptómicos en melanoma a partir de cohortes públicas. La principal decisión metodológica ha sido separar el análisis predictivo en descubrimiento, validación y sensibilidad, evitando mezclar tejidos, plataformas y checkpoints en un único análisis principal.

Los resultados iniciales identifican una señal biológicamente coherente relacionada con presentación antigénica MHC-I y respuesta a IFN-gamma. Sin embargo, el tamaño muestral y la heterogeneidad de plataformas obligan a interpretar los hallazgos como exploratorios.

\section{Líneas de futuro}

Las siguientes líneas de trabajo consisten en consolidar el bloque pronóstico TCGA-SKCM, completar la discusión con bibliografía específica, generar figuras finales para la memoria y evaluar si procede incorporar modelos predictivos regularizados, como lasso o elastic net, manteniendo una validación interna estricta.

\section{Seguimiento de la planificación}

La planificación inicial se ha modificado para responder a limitaciones metodológicas detectadas durante el desarrollo. El cambio principal ha sido reinterpretar el dataset integrado GEO como análisis exploratorio y priorizar el descubrimiento en GSE78220. Este cambio se considera necesario y adecuado, ya que mejora la coherencia biológica del diseño y responde a las recomendaciones de la tutora.

\section{Evaluación de impactos ético-sociales, sostenibilidad y diversidad}

Los impactos previstos en la sección \ref{s:etic} se han abordado mediante uso de datos públicos anonimizados, documentación reproducible, control de versiones y declaración explícita de limitaciones. La perspectiva de género se considera en TCGA mediante disponibilidad de sexo como covariable potencial, y se documenta la limitación de las cohortes GEO al no disponer de sexo armonizado.
